\chapter{Working with Text}\label{ch:strings}

\begin{goals}
\begin{itemize}
  \item Joining strings and building composite messages
  \item Special characters: new line, tab and a quotation mark inside a string
  \item Cutting a piece out of a string, splitting it into words and trimming spaces
  \item Searching inside a string and comparing strings
  \item Converting strings to numbers and back
\end{itemize}
\end{goals}

\section{Text is everywhere}

Names, messages, addresses, phone numbers, dates: a large part of the data
in any program is text. We have worked with strings in the earlier chapters,
but always as one solid piece. In this chapter we learn to look inside a
string, take a piece out of it, split it apart and draw information from
it. All of this is done with \textbf{capabilities} that are attached to a
string with a dot, like the \code{to\_int} you saw in Chapter~\ref{ch:types}.

\section{Joining and length}

\program{Joining strings}
\begin{salamcode}
func main:
    first name := "Sara"
    last name := "Ahmadi"
    full name := first name + " " + last name
    println full name
    phone := "09121234567"
    println "Length of the number:", phone.len()
end
\end{salamcode}

\begin{salamout}
Sara Ahmadi
Length of the number: 11
\end{salamout}

\code{len()} gives the number of compartments the string takes up in
memory. For a phone number, or any text made of digits and English letters,
each character is one compartment, and this number is what you would
expect.

\begin{warn}[The length of text in other alphabets]
The computer keeps every English letter, digit and space in \emph{one}
compartment, but letters with accents such as \code{é}, and letters of other
alphabets such as Persian, Arabic or Greek, take two or more compartments
each. \code{len()} counts compartments, so \code{"café".len()} is 5, not 4,
and \code{"γεια".len()} is 8, not 4. For this reason, in this book we
use \code{len()} and \code{substr()} mostly on text made of digits and
English letters (numbers, dates, codes). Salam has other tools for counting
such letters exactly, but they are outside the scope of this book.
\end{warn}

\section{Special characters}

Some things cannot be written directly inside quotation marks: moving to the
next line, or the quotation mark itself. For these Salam has combinations
that begin with a backslash \code{\textbackslash}:

\begin{center}
\begin{tabular}{cl}
\toprule
Written as & Meaning \\
\midrule
\code{\textbackslash n} & move to the next line \\
\code{\textbackslash t} & a tab (one long jump) \\
\code{\textbackslash "} & the quotation mark itself \\
\code{\textbackslash\textbackslash} & the backslash itself \\
\bottomrule
\end{tabular}
\end{center}

\program{Special characters}
\begin{salamcode}
func main:
    println "first line\nsecond line"
    println "Name: \t Sara"
    println "She said: \"hello\""
end
\end{salamcode}

\begin{salamout}
first line
second line
Name:      Sara
She said: "hello"
\end{salamout}

With \code{\textbackslash n}, one \code{println} can write several lines.
Each of these combinations counts as \emph{one} character, even though it is
written with two symbols.

\section{Other kinds of quotation marks}

In this book we always use \code{"..."}; learn that, and it is enough. But
Salam has three other markers for writing fixed text in code which you may
see in other people's programs, so it is worth knowing them:

\begin{center}
\begin{tabular}{c>{\raggedright\arraybackslash}p{5.6cm}c}
\toprule
Marker & Meaning & Escapes \\
\midrule
\code{"..."} & a string of any length & yes \\
\code{'...'} & exactly \emph{one} character (type \code{char}), not a string & yes \\
\code{\textasciigrave...\textasciigrave} & a raw string; no sequence is processed & no \\
\code{u'...'} & exactly one Unicode character from any alphabet (type \code{uchar}) & no \\
\bottomrule
\end{tabular}
\end{center}

\program{Three other kinds}
\begin{salamcode}
func main:
    letter := 'A'
    println letter
    accented := u'é'
    println accented
    path := `C:\Users\Sara`
    println path
end
\end{salamcode}

\begin{salamout}
A
é
C:\Users\Sara
\end{salamout}

Look at the third line of the output: in a raw string between
\code{\textasciigrave} and \code{\textasciigrave}, the backslash no longer
starts an escape sequence and is printed just as written; this is very handy
for writing a Windows path or a search pattern full of
\code{\textbackslash}.

The first two lines show the two character types. \code{'A'} has the type
\code{char}, a single character that fits in one byte, and \code{u'é'} has
the type \code{uchar}, a single character from any alphabet. A
\code{char} is stored as a number, so you can convert it with \code{as}:
\code{'A' as int} is 65, the code the computer uses for the letter A, and
\code{65 as char} is \code{A} again. A \code{uchar} converts in the same
way, to and from its Unicode code point: \code{u'é' as int} is 233, and
\code{1587 as uchar} is the Persian letter \fa{س}.

\begin{warn}[\code{'...'} holds one byte only]
\code{'...'} accepts only \emph{one byte}, not a whole character. English
letters are one byte each, so \code{'A'} is fine. But a letter such as
\code{é}, and every letter of alphabets such as Persian or Greek, takes two
bytes or more, so \code{'é'} meets this error: ``a char literal holds one
byte; write \code{u'...'} for a Unicode character''. When you need one
Unicode character (from any alphabet), use \code{u'...'} instead of
\code{'...'}, as in the fourth line of the program above.
\end{warn}

\section{Cutting a piece out of a string}

\code{substr(start, count)} takes \code{count} characters beginning at
position \code{start} (which, like array indices, is counted from zero):

\program{The pieces of a date}
\begin{salamcode}
func main:
    date := "2025/09/19"
    year := date.substr(0, 4)
    month := date.substr(5, 2)
    day := date.substr(8, 2)
    println "Year:", year, "Month:", month, "Day:", day
    println "Next year:", year.to_int() + 1
end
\end{salamcode}

\begin{salamout}
Year: 2025 Month: 09 Day: 19
Next year: 2026
\end{salamout}

Count the positions: the four digits of the year are at positions 0 to 3,
then a slash at position 4, then the month at 5 and 6, a slash at 7, and the
day at 8 and 9. The result of \code{substr} is a new string; the original
string stays untouched.

With the same tool we can check a phone number:

\program{Checking a mobile number}
\begin{salamcode}
func main:
    phone := "09121234567"
    prefix := phone.substr(0, 4)
    if (phone.len() == 11) and (prefix == "0912"):
        println "This is a mobile number on the first network"
    else:
        println "Unknown number"
    end
end
\end{salamcode}

\begin{salamout}
This is a mobile number on the first network
\end{salamout}

\section{Splitting into words}

\code{split(separator)} cuts the string wherever the separator appears and
gives a list of the pieces. This list is like an array: we read its elements
with \code{get(index)} (or with \code{[index]}, as with an array) and we
can walk through it with \code{each}:

\program{Counting the words of a sentence}
\begin{salamcode}
func main:
    sentence := "Salam is a language for everyone"
    words := sentence.split(" ")
    println "Number of words:", words.len()
    repeat words.len() in i:
        println i + 1, "-", words.get(i)
    end
end
\end{salamcode}

\begin{salamout}
Number of words: 6
1 - Salam
2 - is
3 - a
4 - language
5 - for
6 - everyone
\end{salamout}

The separator can be any string: \code{split(",")} for a list separated by
commas, or \code{split("/")} for the date above.

\section{Trimming extra spaces}

Text that comes from a user or from a file often has extra spaces at the
beginning and the end. \code{trim()} removes them:

\program{Trimming}
\begin{salamcode}
func main:
    raw text := "   pass   "
    clean := raw text.trim()
    println "[" + raw text + "]"
    println "[" + clean + "]"
    println clean == "pass"
end
\end{salamcode}

\begin{salamout}
[   pass   ]
[pass]
true
\end{salamout}

We put the square brackets in so that the invisible spaces can be seen.
Without trimming, the comparison with \code{"pass"} would have been false,
because the spaces are part of the string too. A good rule: trim every piece
of text that comes from outside before you compare it.

\section{Searching inside a string}

\code{find(text)} says where the given text begins inside the string, and
if it is not in the string at all, it returns \code{-1}. Its most common
use is exactly this question of ``is it there or not'':

\program{Searching a message}
\begin{salamcode}
func main:
    message := "payment completed successfully"
    if message.find("success") != -1:
        println "The operation succeeded"
    else:
        println "The operation failed"
    end
    println message.find("error")
end
\end{salamcode}

\begin{salamout}
The operation succeeded
-1
\end{salamout}

\section{Strings and numbers}

In Chapter~\ref{ch:types} we saw that \code{+} between a string and a number
converts the number to a string and joins them, and that \code{to\_int()}
and \code{to\_float()} go the other way:

\program{From text to a number}
\begin{salamcode}
func main:
    price text := "25"
    count text := "3"
    price := price text.to_int()
    count := count text.to_int()
    println "Amount:", price * count
    weight := "2.5".to_float()
    println "Weight of two crates:", weight * 2.0
end
\end{salamcode}

\begin{salamout}
Amount: 75
Weight of two crates: 5
\end{salamout}

\begin{warn}
If the string is not a number at all (such as \code{"abc"}),
\code{to\_int()} and \code{to\_float()} give zero, so make sure the text
really is a number before you convert it.
\end{warn}

\section{Building a string with a loop}

In Chapter~\ref{ch:loops} we built a string piece by piece (the triangle of
stars). That job can be wrapped up in a function that repeats any string any
number of times:

\program{Repeating a string}
\begin{salamcode}
func times(text: str, count: int): str:
    mut result := ""
    repeat count:
        result += text
    end
    ret result
end

func main:
    println times("-=", 10)
    println times("*", 5) + " five stars " + times("*", 5)
    println times("-=", 10)
end
\end{salamcode}

\begin{salamout}
-=-=-=-=-=-=-=-=-=-=
***** five stars *****
-=-=-=-=-=-=-=-=-=-=
\end{salamout}

\section{A more complete program}

We take a date in the form \code{"2025/09/19"} and print it in a readable
form, with the name of the month:

\program{A readable date}
\begin{salamcode}
func month name(number: int): str:
    ret match number:
        1: "January" end
        2: "February" end
        3: "March" end
        4: "April" end
        5: "May" end
        6: "June" end
        7: "July" end
        8: "August" end
        9: "September" end
        10: "October" end
        11: "November" end
        12: "December" end
        else: "invalid" end
    end
end

func main:
    date := "2025/09/19"
    year := date.substr(0, 4)
    month := date.substr(5, 2).to_int()
    day := date.substr(8, 2).to_int()
    println day, month name(month), year
end
\end{salamcode}

\begin{salamout}
19 September 2025
\end{salamout}

Look at the line \code{date.substr(5, 2).to\_int()}: the result of
\code{substr} is a string, and another capability can be called on it
straight away. This ``chaining'' is very common in work with strings.

\begin{summary}
\begin{itemize}
  \item \code{+} joins strings and \code{len()} gives the number of
        characters (accented and non-Latin letters count as two or more).
  \item \code{\textbackslash n}, \code{\textbackslash t} and
        \code{\textbackslash "} are special characters.
  \item \code{substr(start, count)} cuts out a piece,
        \code{split(separator)} splits a string into a list of pieces,
        \code{trim()} removes the spaces at both ends, and
        \code{find(text)} gives the position of a text (or \code{-1}).
  \item \code{to\_int()} and \code{to\_float()} convert a string to a
        number.
  \item Capabilities can be chained: \code{text.substr(0, 4).to\_int()}.
  \item A character converts to its code and back with \code{as}:
        \code{'A' as int} is 65 and \code{65 as char} is \code{A};
        \code{u'é' as int} is 233 and \code{233 as uchar} is \code{é}.
\end{itemize}
\end{summary}

\section*{Exercises}
\markright{Exercises}

\exercise Put a first name and a last name into two variables and, with a
single \code{println}, print three lines: the full name, the last name on
its own, and a sentence that contains both. (Hint: \code{\textbackslash n}.)

\exercise Put a ten-digit identity number into a string and check that it
has exactly ten characters. Then separate its first three digits and print
them.

\exercise Put a sentence into a variable and find and print its longest
word.

\exercise Put a list of grades into a variable as one comma-separated string
(\code{"17,15,19,12"}). Split it, convert each piece to a number, and print
the sum and the average.

\exercise Write a function that takes a string and a width and, if the
string is shorter than the width, attaches enough spaces to its right to
bring it up to that width. Use it to print a neat two-column table of
product names and prices.

\exercise Print, on one line, the five letters that come after \code{u'\fa{س}'}
in the Unicode table. Convert the letter to its code with \code{as int}, add
1 to 5 to it, and convert each result back with \code{as uchar}.
